\documentclass{ximera}

\title{Derivative Review Activity Faciliatator Guide}
\author{MATH 426: Calculus II}

\begin{document}
\begin{abstract}
   Working with the peers in your group, solve the following problems. Make sure to show and justify all your work. Make sure everyone in the group understands the solution and participates. Be prepared to report your answers to the whole class. 
\end{abstract}
\maketitle

\textbf{Student Learning Outcomes}\\
Students will be able to:
    \begin{itemize}
        \item Recall derivatives of known functions.
        \item Recall derivative rules.
        \item Discuss and brainstorm rules from Calculus I or before that may be helpful in Calculus II.
    \end{itemize}

In large part, Calculus I was the study of differentiation.  In a similar way, Calculus II is a study of integration (for at least part of the semester).  This activity is meant to be review, but also give you a reference sheet which you can look back at during Calculus II: we are creating a reference sheet and study guide.  Keep this document in your notebook or somewhere else where it will not get lost.

\begin{exercise}
   \textbf{Known Derivatives}
   \begin{itemize}
    \item $\frac{d}{dx}(c)=$ \textcolor{blue}{$0$} (where $c$ is any constant.) 
    \item $\frac{d}{dx}(x^n)=$ \textcolor{blue}{$nx^{n-1}$}
    \item $\frac{d}{dx}(e^x)=$ \textcolor{blue}{$e^x$}
    \item $\frac{d}{dx}(\ln|x|)=$ \textcolor{blue}{$\frac{1}{x}$}
    \item $\frac{d}{dx}(b^x)=$ \textcolor{blue}{$b^x\ln(b)$} (where $b>0$) 
    \item $\frac{d}{dx}(\log_b(x))=$ \textcolor{blue}{$\frac{1}{x\ln(b)$}} (where $b>0$)
    \item $\frac{d}{dx}(\sin(x))=$ \textcolor{blue}{$cos(x)$}
    \item $\frac{d}{dx}(\cos(x))=$ \textcolor{blue}{$-\sin(x)$}
    \item $\frac{d}{dx}(\tan(x))=$ \textcolor{blue}{$\sec^2(x)$}
    \item $\frac{d}{dx}(\sec(x))=$ \textcolor{blue}{$\tan(x)\sec(x)$}
    \item $\frac{d}{dx}(\csc(x))=$ \textcolor{blue}{$-\cot(x)csc(x)$}
    \item $\frac{d}{dx}(\cot(x))=$ \textcolor{blue}{$-\csc^2(x)$}
    \item $\frac{d}{dx}(\arcsin(x))=$ \textcolor{blue}{$\frac{1}{\sqrt{1-x^2}}$}
    \item $\frac{d}{dx}(\arccos(x))=$ \textcolor{blue}{$-\frac{1}{\sqrt{1-x^2}}$}
    \item $\frac{d}{dx}(\arctan(x))=$ \textcolor{blue}{$\frac{1}{1+x^2}$}
    \item $\frac{d}{dx}(arcsec(x))=$ \textcolor{blue}{$-\frac{1}{1+x^2}$}
    \item $\frac{d}{dx}(arccsc(x))=$ \textcolor{blue}{$\frac{1}{|x|\sqrt{1-x^2}}$}
    \item $\frac{d}{dx}(arccot(x))=$ \textcolor{blue}{$-\frac{1}{|x|\sqrt{1-x^2}}$}
   \end{itemize}
\end{exercise}

\begin{exercise}
    \textbf{Derivative Rules} \\
    When writing the derivative rules, pay close attention to your notation.  You can write it in a way that makes sense to you, but make sure you aren't losing any important mathematical information that might cause you problems later.
    \begin{itemize}
        \item Power rule: \textcolor{blue}{$\frac{d}{dx}(x^n)=nx^{n-1}$}
        \item Sum rule: \textcolor{blue}{$\frac{d}{dx}(f(x)+g(x))=\frac{d}{dx}(f(x))+\frac{d}{dx}(g(x))$}
        \item Difference rule: \textcolor{blue}{$\frac{d}{dx}(f(x)-g(x))=\frac{d}{dx}(f(x))-\frac{d}{dx}(g(x))$}
        \item Constant multiple rule: \textcolor{blue}{$\frac{d}{dx}(cf(x))=c\frac{d}{dx}(f(x))$}
        \item Product rule: \textcolor{blue}{$\frac{d}{dx}(f(x)g(x))=f(x)\frac{d}{dx}(g(x))+g(x)\frac{d}{dx}(f(x))$}
        \item Quotient rule: \textcolor{blue}{$\frac{d}{dx}\left(\frac{f(x)}{g(x)}\right)=\frac{g(x)\frac{d}{dx}(f(x))+f(x)\frac{d}{dx}(g(x))}{(g(x))^2}$}
        \item Chain rule: \textcolor{blue}{$\frac{d}{dx}(g(f(x)))=g'(f(x))f'(x)$}
    \end{itemize}
\end{exercise}

\begin{exercise}
    \textbf{Other important facts:}\\
    What other facts from Calculus I (or before) do you think you might want to remember?  (If you are stuck, this is a great time to ask your TA or LA for their suggestions!)\\
    \textcolor{blue}{This is subjective and unique to your group!  I would suggest that facts about limits at infinity might be helpful for later in Calculus II. MAC tutors and LAs also suggested that trig functions, the Pythagorean trig identites, and log and exponent rules will be helpful!}
\end{exercise}

\end{document}